% ATTEMPT_STATUS: unresolved
% ATTEMPT_ORIGIN: new research, 2026-10-01
% TOP_PROBLEM: {"id": 1394, "title": "Perfect 1-Factorization Conjecture", "category_id": 3, "category": {"id": 3, "name": "graph_theory", "display_name": "Graph Theory", "description": "Problems involving graphs, networks, and their properties.", "slug": "graph-theory", "order_index": 3, "created_at": "2026-07-31T15:26:25.671Z"}, "status": "open"}
% FIRST_POSTED: 2026-10-01
% Imported from: unsolvedmath_additional_500.tex; UnsolvedMath ID 1394
% !TeX program = lualatex
% Compile twice: lualatex 1394.tex
% Self-contained: no external bibliography, images, or \input files.
% Numeric section labels and visible numbers are original UnsolvedMath IDs.
\documentclass[11pt,a4paper]{article}
\usepackage[margin=24mm,headheight=14pt]{geometry}
\usepackage{fontspec}
\setmainfont{Latin Modern Roman}
\setsansfont{Latin Modern Sans}
\setmonofont{Latin Modern Mono}
\usepackage{amsmath,amssymb,mathtools}
\usepackage{unicode-math}
\setmathfont{Latin Modern Math}
\usepackage{microtype}
\usepackage{enumitem,xurl,needspace}
\usepackage[dvipsnames]{xcolor}
\usepackage{hyperref}
\hypersetup{unicode=true,colorlinks=true,linkcolor=MidnightBlue,urlcolor=MidnightBlue,
 pdftitle={UnsolvedMath Problem 1394},pdfauthor={Source-annotated compilation},bookmarksnumbered=true}
\usepackage{bookmark,tocloft,titlesec,fancyhdr}
\setlength{\cftsecnumwidth}{4.2em}
\cftsetpnumwidth{2.5em}
\cftsetrmarg{4em}
\titleformat{\section}{\Large\bfseries\raggedright}{\thesection}{0.7em}{}
\pagestyle{fancy}\fancyhf{}
\fancyhead[L]{\small\sffamily UnsolvedMath research attempt}
\fancyhead[R]{\small Original catalogue IDs}
\fancyfoot[C]{\thepage}
\setlength{\parindent}{0pt}
\setlength{\parskip}{5pt plus 1pt}
\setlength{\emergencystretch}{3em}
\setcounter{tocdepth}{1}\setcounter{secnumdepth}{1}
\setlist[itemize]{leftmargin=3em,itemsep=4pt,topsep=4pt,parsep=0pt}
\allowdisplaybreaks
\newcommand{\N}{\mathbb{N}}
\newcommand{\Z}{\mathbb{Z}}
\newcommand{\Q}{\mathbb{Q}}
\newcommand{\R}{\mathbb{R}}
\newcommand{\C}{\mathbb{C}}
\newcommand{\F}{\mathbb{F}}
\newcommand{\A}{\mathbb{A}}
\newcommand{\PP}{\mathbb{P}}
\newcommand{\E}{\mathbb{E}}
\newcommand{\eps}{\varepsilon}
\newcommand{\dd}{\,\mathrm d}
\newcommand{\sref}[2]{\hyperlink{source:#1:#2}{[S#2]}}
\newif\ifshowcatalogue
\showcataloguetrue
\newcommand{\cataloguescope}[1]{\ifshowcatalogue
 \paragraph{Statement recorded in the supplied source.}
 #1\fi}
\begin{document}
% UnsolvedMath ID 1394; source catalogue code GRAPH-007

\setcounter{section}{1393}

\section{Perfect One-Factorizations of Complete Graphs}\label{1394}

{\small\sffamily UnsolvedMath ID 1394 · Graph Theory}

\subsection{Definitions and mathematical statement}

\paragraph{Shared notation.}Unless a section says otherwise, $\N=\{1,2,\ldots\}$,
$\Z$ is the ring of integers, $\Q,\R,\C$ are the rational, real, and complex
fields, $[N]=\{1,\ldots,\lfloor N\rfloor\}$, and $\log$ is the natural
logarithm. For real functions of a parameter tending to infinity,
$f\ll g$ and $f=O(g)$ mean $|f|\leq Cg$ eventually for some fixed $C$;
$f\gg g$ reverses this comparison; $f=o(g)$ means $f/g\to0$;
$f\sim g$ means $f/g\to1$; and $f\asymp g$ means both order bounds.
A subscript on $O$ or $\ll$ specifies allowed constant dependence.
The expression $n^{o(1)}$ in an upper bound means growth smaller than
$n^\epsilon$ for each fixed $\epsilon>0$. Local definitions govern any
exception, finite-field convention, or use of the same symbol for another
object. An asymptotic claim is not automatically an effective algorithm.



A one-factor is a perfect matching. A one-factorization of $K_{2n}$ partitions its edges into $2n-1$ perfect matchings. It is perfect if the union of any two different matching classes is one Hamiltonian cycle, rather than a disjoint union of shorter even cycles. The conjecture quantifies over all $n\geq2$. \sref{1394}{1} \sref{1394}{2}

\paragraph{Statement recorded in the supplied source.}
Does every complete graph on an even number of vertices admit a perfect 1-factorization? \sref{1394}{1}

\paragraph{Residual task reported by the archive.}

The embedded review (2026-08-17) records: Establish a perfect 1-factorization of every even complete graph. \sref{1394}{1}

\subsection{Short English statement}

Can every even-order complete graph be decomposed into perfect matchings so that any two matchings together form one spanning cycle?

\subsection{Research attempt}
\textbf{Status: UNRESOLVED.} We derive the standard cyclic perfect one-factorization for $K_{p+1}$ with $p$ an odd prime, prove its Hamiltonian property, and give an explicit composite-order failure of the same formula. This infinite special family does not cover every even order.

\paragraph{Source check and target distinction.}
Cheng--Sgueglia [S2], consulted 1 October 2026, proves that a one-factorization can have a subfamily of $(1-o(1))n$ matchings whose pairs are Hamiltonian. Its title's asymptotic wording does not assert that every one of the $n-1$ factors participates in such a perfect family. We keep the exact target unchanged. The attack here is algebraic: construct factors by reflections and analyze the components of their pairwise unions.

\paragraph{A one-factorization for every odd modulus.}
Let $q\ge3$ be odd and take vertex set $\mathbb Z/q\mathbb Z\cup\{\infty\}$. For each $a\in\mathbb Z/q\mathbb Z$ let
\[
F_a=\{\{\infty,a\}\}\cup
\{\{x,2a-x\}:x\ne a\}.
\]
The second collection is a set of unordered pairs, so listing both endpoints does not create repeated edges. The map $x\mapsto2a-x$ is an involution whose only fixed point is $a$, because two is invertible modulo odd $q$. Its remaining vertices form $(q-1)/2$ disjoint pairs. Together with $\{\infty,a\}$ they give a perfect matching.

Every edge involving infinity occurs exactly once. A finite edge $\{x,y\}$ with $x\ne y$ occurs precisely in the factor indexed by $a=(x+y)/2$. Thus the $q$ matchings partition all edges of $K_{q+1}$. This part of the construction needs only oddness, not primality.

\paragraph{Why prime modulus makes every pair Hamiltonian.}
Fix distinct $a,b$ and let $q=p$ be prime. The graph $F_a\cup F_b$ is a simple two-regular spanning graph, hence a disjoint union of cycles; factors share no edge. If it is not connected, some component $C$ avoids infinity. It then also avoids $a$ and $b$, because their respective matching edges lead to infinity.

For every $x\in C$, both reflections $r_a(x)=2a-x$ and $r_b(x)=2b-x$ belong to $C$. Their composition is translation by $2(b-a)$ or its negative, depending on order. This translation is nonzero in $\mathbb F_p$ and has one orbit, the whole field. Closure of a nonempty $C$ under the composition therefore forces $C=\mathbb F_p$. That contradicts $a\notin C$. Every component must consequently contain infinity, so there is just one component, a Hamiltonian cycle.

The proof also applies for an odd composite modulus to a particular pair whenever $\gcd(b-a,q)=1$. It is the transitivity of translation, not a generic appeal to symmetry, that supplies connectedness.

\paragraph{The composite obstruction is real.}
For $q=9$, take $a=0$ and $b=3$. The component containing infinity is the four-cycle
\[
\infty,0,6,3,\infty.
\]
Indeed $\infty0$ and $63$ are in $F_0$, while $06$ and $3\infty$ are in $F_3$. The other six vertices form another component. Thus this one-factorization of $K_{10}$ is not perfect. This is not a counterexample to the existence conjecture, which permits completely different factors. It rules out the unqualified extension of the prime proof to every odd modulus.

For a composite odd $q$, choose a nonzero difference sharing a proper divisor with $q$. The translation orbits split into residue classes, and the argument forcing a single component is no longer available. Changing only the names of the vertices or factors preserves component lengths, so relabeling the failed construction cannot repair it.

\paragraph{Executed checks and certificate format.}
The audit script constructed all factors for $q=3,5,7,9,11,13,17$, verified that each is a perfect matching, checked that their edge sets partition $K_{q+1}$, and counted components for every pair. It found no disconnected pair at the listed prime moduli. At $q=9$ it found nine bad pairs, including $(0,3)$ with a four-vertex component containing infinity. These exact set and graph traversals are recorded in \texttt{root/checks.py} and \texttt{root/checks\_results.json} in the local audit.

For any proposed perfect one-factorization, verifying the target is finite and direct: check the edge partition and traverse every two-factor union. Finding a factorization is the difficult part. A family of many pairwise Hamiltonian factors does not automatically extend by one more factor: the unused edges must support a perfect matching satisfying a connectedness condition against every already chosen factor simultaneously.

\paragraph{Remaining gap.}
The unresolved task is construction at every even order beyond the prime-successor family, without retaining any non-Hamiltonian pair. The precise gap in this approach is replacing the transitive translation argument when the modulus has nonunits. Local switches in a pair of factors can merge its cycles but can also damage the unions with other factors; no invariant guaranteeing simultaneous improvement is proved here. The result is a checked classical-family derivation and obstruction analysis, with no claimed new solution.

\subsection{Sources}

{\small

\begin{itemize}

\item[\hypertarget{source:1394:1}{S1}] ULAM.AI, \emph{UnsolvedMath}, supplied \texttt{problems.json}, record 1394 (GRAPH-007), statement and background. The embedded literature-review date is 2026-08-17; that is the archive’s review date, not a fresh certification. Dataset landing page: \url{https://huggingface.co/datasets/ulamai/UnsolvedMath}.

\item[\hypertarget{source:1394:2}{S2}] Y. Cheng and A. Sgueglia, The perfect 1-factorisation conjecture holds asymptotically, arXiv:2607.09459 (2026). \url{https://arxiv.org/abs/2607.09459}. Bibliographic information imported from the supplied record.

\end{itemize}

}

\label{end:1394}
\end{document}
